% SHARD-ID: CA-32-GAUSSIAN-HIGHER-DIMENSIONAL-CELL-CURRENTS
% SHARD-TITLE: Higher-Dimensional Gaussian Cell-Current Proposal
% SHARD-SUMMARY: Proposes a compiler-input specification for d=2,3 Gaussian cell currents using oriented cells, faces, and incidence signs.
% SHARD-SUMMARY: Separates energy flux, momentum density, and stress witnesses instead of inferring T_a0 from T_0a or assuming symmetric stress.
% SHARD-SUMMARY: Records rotation, trace, improvement, and open-versus-periodic boundary diagnostics as future numerical targets, not convergence claims.
% SHARD-KEYWORDS: Gaussian boson, higher dimensions, cell current, stress tensor, rotation, boundary terms, compiler input

\section{Higher-Dimensional Gaussian Cell-Current Proposal}
\label{sec:gaussian-higher-dimensional-cell-currents}

\subsection{Status, Sources, and Dependencies}

\textbf{Status: proposal and compiler-input specification.}  This shard states
the support data and local equations a higher-dimensional Gaussian
stress-energy candidate must carry in \(d=2,3\).  It is not a theorem, does not
define a global orientation convention, and makes no convergence claim.

Dependencies: CA-13, CA-23, CA-29--CA-31, and CONVENTIONS.md (j), (k), (l).

% Source: references/qft/SOURCES.md:59--90 -- registered Weinberg A1 anchor ranges and improvement-status boundary.
% Source: references/qft/SOURCES.md:94--97 -- targeted OCR is an anchor convenience; exact signs/indices need page-image checks.
% Source: references/qft/Weinberg1995/Weinberg1995_QFT1_A1_ocr.txt:260--280 -- translations give conserved currents grouped as the energy-momentum tensor and integrated translation generators.
% Source: references/qft/Weinberg1995/Weinberg1995_QFT1_A1_ocr.txt:357--361 -- canonical raised tensor need not be symmetric.
% Source: references/qft/Weinberg1995/Weinberg1995_QFT1_A1_ocr.txt:572--589 -- Belinfante tensor is conserved, has the same integrated translations, and is symmetric.
% Source: references/qft/Weinberg1995/Weinberg1995_QFT1_A1_ocr.txt:590--633 -- Lorentz generators from the symmetric tensor, including rotations and boosts.
% Source: references/qft/Weinberg1995/Weinberg1995_QFT1_A1_ocr.txt:675--684 -- Lorentz commutator proof caveat for boost generators with auxiliary fields.
% Convention: CONVENTIONS.md:196--239 -- finite periodic Gaussian checks do not define periodic first moments or finite coordinate generators.
% Convention: CONVENTIONS.md:260--277 -- flat-symbol boost-time target and non-identification with sourced OAR or mass-shell representations.
% Convention: CONVENTIONS.md:312--347 -- Poincare vector fields translated to self-adjoint commutator signs.
% Convention: CONVENTIONS.md:349--421 -- Gaussian density split, boundary, subtraction, and improvement policies.
% Source: report/sections/13_position_dependent_bulk_generators.tex:65--94 -- boosts start from energy density, while rotations require momentum-density witnesses.
% Source: report/sections/13_position_dependent_bulk_generators.tex:96--105 -- compiler needs geometry and no orientation/convergence convention is fixed there.
% Source: report/sections/29_gaussian_boson_real_space_density.tex:26--42 -- physical density h_x versus integrated cell energy e_x.
% Source: report/sections/29_gaussian_boson_real_space_density.tex:93--105 -- open-boundary and periodic first-moment policies.
% Source: report/sections/29_gaussian_boson_real_space_density.tex:122--151 -- improvement terms shift first moments by edge and boundary terms.
% Source: report/sections/30_gaussian_boson_1d_stress_candidates.tex:83--146 -- T_00, T_01, T_10, and T_11 candidate roles in one dimension.
% Source: report/sections/30_gaussian_boson_1d_stress_candidates.tex:184--194 -- pending momentum-density, stress-current, and continuum witnesses.
% Source: report/sections/31_gaussian_current_symbol_equivalence.tex:55--98 -- one-dimensional current-symbol equality and its limited status.
% Source: report/sections/31_gaussian_current_symbol_equivalence.tex:122--137 -- finite-chain momentum, periodic first moments, and convergence remain open.

\subsection{Compiler Input: Oriented Support Data}

For each finite or infinite hypercubic support in \(d=2\) or \(d=3\), the
compiler input must include an oriented incidence structure:
\[
\mathcal C,\qquad \mathcal F,\qquad
\sigma(c,f)\in\{-1,0,+1\}.
\]
Here \(\mathcal C\) is the set of energy cells, attached to sites \(x_c\in
\epsilon\mathbb Z^d\).  The set \(\mathcal F\) contains oriented codimension-one
faces.  In \(d=2\), these faces are oriented edges between neighbouring square
cells.  In \(d=3\), they are oriented plaquettes between neighbouring cube
cells.  If an interior face \(f\) is oriented from its minus cell \(c_-(f)\) to
its plus cell \(c_+(f)\), set
\[
  \sigma(c_-(f),f)=+1,\qquad
  \sigma(c_+(f),f)=-1,
\]
and set \(\sigma(c,f)=0\) for all other cells.  This sign convention is local
input for this proposal, not a repository-wide orientation convention.

Boundary faces are separate data.  For an open box, \(\mathcal F_\partial\)
contains oriented exterior faces with only one adjacent cell; their incidence
sign is recorded as \(\sigma_\partial(c,f)=+1\) for outward flux.  Setting a
boundary flux to zero is then a named boundary condition, not an omitted term.
For a periodic support, exterior faces are absent, but any first moment or
coordinate-weighted test still needs the branch or Fourier policy required by
CONVENTIONS.md (j).

\subsection{Energy Density and Energy Flux}

The energy component starts from CA-29:
\[
  T_{00}^{\mathrm{lat}}(c):=h_c,\qquad
  e_c:=\epsilon^d h_c,\qquad
  H_\epsilon=\sum_{c\in\mathcal C}e_c .
\]
All finite commutator identities use the integrated cell energy \(e_c\).

An energy-current candidate is face data
\[
  J_0(f)\quad (f\in\mathcal F),
\]
with optional boundary data \(J^\partial_0(f)\) on \(f\in\mathcal F_\partial\).
If \(f\) has normal direction \(a\), then \(J_0(f)\) is the proposed lattice
\(T_{0a}\) flux through that oriented face.  The required energy-continuity
witness is the cell equation
\[
  i[H_\epsilon,e_c]
  +
  \sum_{f\in\mathcal F}\sigma(c,f)J_0(f)
  +
  \sum_{f\in\mathcal F_\partial}\sigma_\partial(c,f)J^\partial_0(f)
  =0 .
\]
For an open no-leak boundary this specializes to \(J^\partial_0(f)=0\).  For a
driven or truncated-region diagnostic, the same displayed equation records the
boundary residual explicitly.  This shard does not propose a universal formula
for \(J_0(f)\); a later numerical suite must derive or solve for it from the
chosen Gaussian density and then test the equation above.

\subsection{Momentum Densities and Stress Currents}

The momentum-density components are additional witnesses:
\[
  M_a(c)\quad (a=1,\ldots,d),
  \qquad
  P_a^{\mathrm{lat}}=\sum_c M_a(c).
\]
They are not inferred from the energy flux \(J_0(f)\).  The reason is already
visible in CA-30: \(T_{10}\) was kept as a momentum-density target, not as a
consequence of \(T_{01}\).  In higher dimensions this separation is mandatory
because rotations require \(M_a\) before angular momentum even has a candidate
density.

For each chosen momentum density \(M_a(c)\), the stress components are its
currents.  If \(f\) has normal direction \(b\), write
\[
  S_a(f)\equiv T_{ab}^{\mathrm{lat}}(f).
\]
The required momentum-continuity witness is
\[
  i[H_\epsilon,M_a(c)]
  +
  \sum_{f\in\mathcal F}\sigma(c,f)S_a(f)
  +
  \sum_{f\in\mathcal F_\partial}\sigma_\partial(c,f)S^\partial_a(f)
  =0
\]
for every \(a\) and every cell \(c\).  Until \(M_a\) and \(S_a\) are supplied
with this witness, there is no lattice \(T_{a0}\) or \(T_{ab}\) claim.

\subsection{Boost and Symbol Diagnostics}

The boost first moment can be formed from the integrated energy cells:
\[
  K_a^{\mathrm{lat}}:=\sum_c x_c^a e_c .
\]
Using the energy-continuity equation, \(i[H_\epsilon,K_a^{\mathrm{lat}}]\)
reduces by discrete summation by parts to the integrated \(a\)-normal energy
flux plus boundary terms.  The higher-dimensional analogue of CA-31 is then a
test, not an assumption:
\[
  \text{integrated }T_{0a}\text{ symbol}
  \quad\stackrel{?}{=}\quad
  \frac12\partial_a\omega(k)^2 .
\]
Only after this equality is checked in the relevant local or symbol regime may
it be compared with the flat one-particle momentum target \(k_a\) and the
residuals of CA-26--CA-27.  Passing the symbol equality would still not prove a
continuum stress tensor.

\subsection{Rotation and Antisymmetric-Stress Residual}

For \(a<b\), a rotation candidate requires the momentum-density witnesses:
\[
  J_{ab}^{\mathrm{lat}}
  :=
  \sum_c \bigl(x_c^aM_b(c)-x_c^bM_a(c)\bigr).
\]
This formula is unavailable if only \(T_{00}\) and \(T_{0a}\) have been supplied.
Applying the momentum-continuity equation gives a discrete summation-by-parts
diagnostic of the form
\[
  i[H_\epsilon,J_{ab}^{\mathrm{lat}}]
  =
  \epsilon\sum_{f\perp a} S_b(f)
  -
  \epsilon\sum_{f\perp b} S_a(f)
  +
  \text{boundary terms},
\]
with the displayed sign tied to the incidence convention above.  In continuum
language this is the antisymmetric-stress obstruction: rotation conservation
targets \(T_{ab}=T_{ba}\), after boundary terms and any spin/improvement data
have been named.  Weinberg's registered anchors motivate that a symmetric
Belinfante tensor can carry the same integrated translations, but the canonical
tensor need not be symmetric.  Therefore stress symmetry is a target or
witness, not an assumption baked into the lattice proposal.

\subsection{Trace, Improvement, and Boundary Alternatives}

Trace diagnostics are also target-dependent.  A compiler may report the
diagonal stress combination built from \(h_c\) and the \(S_a(f)\), but it must
name the continuum target and metric/sign convention before calling it a trace
residual.  No massless-CFT tracelessness test, massive Klein-Gordon trace test,
or improvement formula is imported here.

The CA-29 improvement policy remains active:
\[
  h'_c=h_c+(\nabla_\epsilon\cdot b)_c .
\]
Any such change must name the edge or face field \(b\), the discrete divergence,
and the boundary convention.  It must also recompute \(J_0\), \(M_a\), \(S_a\),
and the first-moment boundary terms; improvements are not silent gauge moves for
boosts or rotations.

Two boundary regimes are admissible compiler inputs:
\[
\begin{array}{ll}
\text{open box} &
  \text{use explicit }\mathcal F_\partial\text{ terms and test endpoint/face
  residuals},\\
\text{periodic box} &
  \text{use no boundary faces, but fix a coordinate branch, sawtooth,}\\
  &\text{or Fourier-interpolation policy before any first-moment test}.
\end{array}
\]
The existing Gaussian numerical layer is periodic-symbol first.  The next
stress-energy numerical suite, S5 in the action plan, should therefore start
with open-boundary local continuity tests or with a separately sourced periodic
coordinate construction, and should include anisotropic/doubler failure
witnesses at the current-density level.
